% =====================================================================
%  1bat-ud2-8.tex  —  HORA 8: Gimnàstica algebraica (mecanització)
%  (sense preàmbul: s'inclou des de main.tex)
% =====================================================================
\horatitol{Sessió 8 — Gimnàstica algebraica}%
{Mecanitzar les operacions amb polinomis (suma, resta i producte) per fer-les
àgils i sense errors de signe ni de grau.}

\subsection{Idea clau}

A les sessions anteriors hem donat sentit real a les operacions amb polinomis:
sumar pressupostos, restar retallades, multiplicar per un percentatge de
subvenció. Avui l'objectiu és que la \textbf{mecànica es torni automàtica}:
fer les operacions de manera ràpida i segura, sense haver de pensar cada cop
en el context. Quan la mecànica és fluïda, l'atenció queda lliure per
interpretar els resultats.

\begin{idea}
\textbf{Tres operacions, tres paranys a evitar:}
\begin{itemize}[leftmargin=1.4em, itemsep=3pt]
  \item \textbf{Suma:} agrupa primer tots els termes semblants.
        $3x^2 + x^2 = 4x^2$, però $3x^2 + x \neq 4x^2$.
  \item \textbf{Resta:} el signe «$-$» davant del parèntesi canvia
        \emph{tots} els signes de dins.
        $(5x+3)-(2x-1) = 5x+3-2x+1 = 3x+4$
        (\textbf{compte:} $-(-1) = +1$).
  \item \textbf{Producte nombre$\times$polinomi:} el nombre multiplica
        \emph{tots} els termes. $3\per(2x^2-x+4) = 6x^2-3x+12$.
        El grau del polinomi no canvia.
\end{itemize}
\end{idea}

\subsection{Exemples}

\begin{exemple}[escalfament — identificar termes semblants i grau]
Classifica els termes del polinomi $Q = 5x^2 - 3x + 2x^2 + 7 - x$ i
simplifica'l.

\medskip
\begin{tabular}{@{}p{0.88\linewidth}@{}}
Termes en $x^2$: $5x^2$ i $2x^2$ → suma: $7x^2$. \\[3pt]
Termes en $x$: $-3x$ i $-x$ → suma: $-4x$. \\[3pt]
Termes independents: $7$ → queda $7$. \\[3pt]
\textbf{Resultat:} $Q = 7x^2 - 4x + 7$ \quad (grau~2).
\end{tabular}
\end{exemple}

\begin{exemple}[operació combinada: suma i producte]
Calcula $2\per A(x) + B(x)$, amb $A(x) = 3x^2 + x - 5$ i
$B(x) = x^2 - 4x + 2$.
\begin{align*}
2\per A(x) + B(x)
  &= 2\per(3x^2+x-5) + (x^2-4x+2) \\
  &= (6x^2+2x-10) + (x^2-4x+2)      &&\text{(distributiva)} \\
  &= 6x^2+x^2 + 2x-4x + (-10+2)     &&\text{(agrupo semblants)} \\
  &= 7x^2 - 2x - 8.                  &&\text{(simplifico)}
\end{align*}
\end{exemple}

\subsection{Exercicis resolts}

\begin{exresolt}[resta amb canvi de tots els signes]
Calcula $(4x^2 - 2x + 9) - (x^2 + 5x - 3)$.
\solucio
\begin{align*}
&(4x^2-2x+9)-(x^2+5x-3) \\
&= 4x^2-2x+9 - x^2-5x+3   &&\text{(canvio tots els signes: $-x^2$, $-5x$, $+3$)} \\
&= (4-1)x^2 + (-2-5)x + (9+3) \\
&= 3x^2 - 7x + 12.
\end{align*}
\resposta{$3x^2 - 7x + 12$}
\end{exresolt}

\begin{exresolt}[operació combinada: producte i resta]
Calcula $3\per P(x) - 2\per Q(x)$, amb $P(x) = x^2 + 2x - 1$ i
$Q(x) = x^2 - x + 4$.
\solucio
\begin{align*}
3\per P(x) - 2\per Q(x)
  &= 3\per(x^2+2x-1) - 2\per(x^2-x+4) \\
  &= (3x^2+6x-3) - (2x^2-2x+8)       &&\text{(distributiva a cada terme)} \\
  &= 3x^2+6x-3 - 2x^2+2x-8           &&\text{(canvio tots els signes)} \\
  &= x^2 + 8x - 11.
\end{align*}
\resposta{$x^2 + 8x - 11$}
\end{exresolt}

\subsection{Exercicis per fer}

\begin{exercicis}
\begin{exlist}
  \item Simplifica agrupant termes semblants:
    \begin{enumerate}[label=\alph*), leftmargin=1.6em, topsep=2pt]
      \item $4x^2 + 3x - x^2 + 5 - 2x$
      \item $7x^3 - 2x^2 + x^3 + 4x^2 - 3x + 1$
    \end{enumerate}

  \item Calcula les sumes i restes:
    \begin{enumerate}[label=\alph*), leftmargin=1.6em, topsep=2pt]
      \item $(6x^2 + x - 4) + (2x^2 - 5x + 7)$
      \item $(3x^2 - 4x + 1) - (x^2 + 2x - 6)$
      \item $(x^2 + x + 1) - (x^2 + x + 1)$
    \end{enumerate}

  \item Calcula els productes:
    \begin{enumerate}[label=\alph*), leftmargin=1.6em, topsep=2pt]
      \item $5\per(2x^2 - x + 3)$
      \item $-2\per(x^2 + 4x - 1)$
      \item $\num{0,25}\per(8x^2 - 4x + 12)$
    \end{enumerate}

  \item Calcula les operacions combinades:
    \begin{enumerate}[label=\alph*), leftmargin=1.6em, topsep=2pt]
      \item $2\per(x^2 + 3x) + (x^2 - x + 5)$
      \item $3\per(2x - 1) - 2\per(x + 4)$
    \end{enumerate}

  \item \textbf{(Compte amb el signe)} Calcula i compara:
    \begin{enumerate}[label=\alph*), leftmargin=1.6em, topsep=2pt]
      \item $(8x^2 - 3x + 6) - (8x^2 - 3x - 4)$
      \item $(8x^2 - 3x + 6) + (-8x^2 + 3x + 4)$
    \end{enumerate}
    Abans de calcular, prediu si els resultats seran iguals. Explica per
    què sí o per què no un cop ho hagis comprovat.

  \item \textbf{(Raonament)} Sense fer cap càlcul, indica quin és el grau
    del resultat de cadascuna de les operacions:
    \begin{enumerate}[label=\alph*), leftmargin=1.6em, topsep=2pt]
      \item $(5x^3 + 2x) + (x^2 - 1)$
      \item $(5x^3 + 2x^2) - (5x^3 - x + 4)$
      \item $7\per(3x^4 - x^2 + 2)$
    \end{enumerate}
    Explica com has determinat el grau en cada cas sense operar.
\end{exlist}
\end{exercicis}

% ---------------------------------------------------------------------
\fullsolucions{Sessió 8}
\begin{solucions}
\begin{sollist}
  \item (a) $3x^2+x+5$ \quad (b) $8x^3+2x^2-3x+1$
  \item (a) $8x^2-4x+3$ \quad
        (b) $2x^2-6x+7$ \quad
        (c) $0$ (el polinomi menys ell mateix és sempre zero)
  \item (a) $10x^2-5x+15$ \quad
        (b) $-2x^2-8x+2$ \quad
        (c) $2x^2-x+3$
  \item (a) $3x^2+5x+5$ \quad (b) $4x-11$
  \item (a) $10$ \quad (b) $10$. Els resultats són iguals perquè restar un
        polinomi equival a sumar el seu oposat: $-(-4)=+4$ en tots dos casos.
  \item (a) grau~3 (el terme $5x^3$ no té cap igual que el cancel·li) \quad
        (b) grau~2 (els termes $5x^3$ es cancel·len; queda $2x^2+\ldots$) \quad
        (c) grau~4 (multiplicar per un nombre no canvia el grau)
\end{sollist}
\end{solucions}
